% !TEX root = main.tex
% TEMPORARY: plan for restructuring Sec. 5 (Power Particles Revisited), for internal discussion only.
\documentclass[main.tex]{subfiles}
\begin{document}

\renewcommand{\thesection}{\Alph{section}}  % the appendix numbering set in appendix.tex does not survive the subfile group
\section{[Temporary] Plan for Section~\ref{sec:flow}}
\label{sec:appendix-flow-plan}

This temporary section outlines a restructuring of \S\ref{sec:flow} around four candidate integrators and two prior methods.
Throughout, $\smash{\tilde{\qq}^{k+1} = \qq^k + \tstep\,\vel^k}$ denotes the free drift of the particles, $\smash{h = \tstep^2/\stiff^2}$, and $\smash{\norm{\cdot}_{\pvol}}$ and $\smash{P_{\achievable}}$ denote the volume-weighted norm and the projection onto $\achievable$ in this norm.

\subsection*{Key observation: one family of position updates}

Setting both the sites and the target centroids of $\jointdual_\alpha$ to the drift $\smash{\tilde{\qq}^{k+1}}$, the pushed targets become $\smash{(1 + 1/\alpha)\,\tilde{\qq}^{k+1}}$, so by \eqref{eq:regularized-projection} the new centroids are
\begin{equation*}
  \qq^{k+1} = P_{\achievable}\bigl( (1 + 1/\alpha)\,\tilde{\qq}^{k+1} \bigr) .
\end{equation*}
Different stiffnesses $\alpha$ then recover the projection-type schemes, including the position step of Power Particles, as listed in Table~\ref{tab:plan-family}.
\begin{table*}[t]
  \caption{[Plan] One family of position updates (sites and targets of $\jointdual_\alpha$ at the drift).}
  \label{tab:plan-family}
  \centering
  \small
  \begin{tabular}{@{}lll@{}}
    \toprule
    $\alpha$ & position update & scheme \\
    \midrule
    $\alpha \to \infty$ & $P_{\achievable}(\tilde{\qq}^{k+1})$ & projection, no potential \\
    $\alpha = 1/h - 1$ & $P_{\achievable}\bigl(\tilde{\qq}^{k+1}/(1 - h)\bigr)$ & implicit, projected incremental potential \\
    $\alpha \to 0$ & centroids of $\arg\max_\weight \dualenergy$ with sites $\tilde{\qq}^{k+1}$ & Power Particles (position step) \\
    \bottomrule
  \end{tabular}
\end{table*}
The last row holds because, with the dilation factor $s = 1 + 1/\alpha$, the limit $\alpha \to 0$ corresponds to $s \to \infty$, and the projection of $s\,\tilde{\qq}^{k+1}$ onto $\achievable$ then tends to the point exposed by the normal $\tilde{\qq}^{k+1}$, \ie the centroids of the volume-constrained power diagram generated by $\tilde{\qq}^{k+1}$ (\S\ref{sec:background}).
This matches the reading of \S\ref{sec:flow-pp}: Power Particles use the drift as a \emph{normal} (sites), whereas the projection uses it as a \emph{target}, and the two differ by one Lloyd step.

The two prior methods thus sit at opposite ends of our formulations:
\begin{itemize}
  \item \citet{Gallouet2018Lagrangian} use $\dualenergy$ as a \emph{potential}, \ie the penalty scheme for $\alpha = 0$;
  \item Power Particles \cite{deGoes2015PowerParticles} use $\dualenergy$ as a \emph{position correction}, \ie the projection scheme for $\alpha \to 0$, with the velocity update replaced by a pressure projection.
\end{itemize}

\subsection*{Proposed flow of Section~\ref{sec:flow}}

\begin{enumerate}
  \item \textbf{Discrete Lagrangian} (\S\ref{sec:flow-lagrangian}): unchanged.
  \item \textbf{Constrained / SHAKE} (\S\ref{sec:flow-centroids}).
    \begin{itemize}
      \item Dual: $\jointdual$ with sites and target centroids at the particles, $\qq = \tcentroid = \qq^k$.
      \item Update: \eqref{eq:flow-shake}, with the magnitude $s$ of the correction along $\dual^k$ found by a 1D search, each trial of which is one projection.
      \item Exact centroids, symplectic, several solves per step.
    \end{itemize}
  \item \textbf{Projection} (\S\ref{sec:flow-projection}).
    \begin{itemize}
      \item Dual: $\jointdual_\alpha$ with sites and target centroids at the drift $\tilde{\qq}^{k+1}$.
      \item Main scheme: implicit, with $\alpha = 1/h - 1$ and $\tstep < \stiff$.
      \item Remark: the explicit-potential variant (current text), which has no time-step limit.
      \item Keep the paragraphs on properties and admissible predictions.
      \item End with the table above, which sets up \S\ref{sec:flow-pp}.
    \end{itemize}
  \item \textbf{Penalty} (\S\ref{sec:flow-softened}).
    \begin{itemize}
      \item Dual: $\jointdual_\alpha$ with sites and target centroids at the particles, used as a potential.
      \item Explicit and symplectic; the centroids lie within an $O(1/\alpha)$ offset of the particles.
      \item \citet{Gallouet2018Lagrangian} is the case $\alpha = 0$, \ie $\dualenergy$.
    \end{itemize}
  \item \textbf{Power Particles} (\S\ref{sec:flow-pp}).
    \begin{itemize}
      \item Dual: $\dualenergy$ with sites at the drift $\tilde{\qq}^{k+1}$, \ie the projection for $\alpha \to 0$.
      \item The velocity update is a pressure projection rather than the position difference.
      \item Keep the current three-paragraph reading, reframed around the table above.
    \end{itemize}
  \item \textbf{Discussion} (\S\ref{sec:flow-discussion}).
    \begin{itemize}
      \item Comparison table of all six schemes (below).
      \item Keep the relation to Arnold's geodesics.
      \item Fold the choice of quadrature into a short remark.
    \end{itemize}
\end{enumerate}

\subsection*{Update rules}

Each scheme as one time step, in terms of the dual it maximizes.

\paragraph{Constrained / SHAKE.}
Maximize $\jointdual$ with $\qq = \tcentroid = \qq^k$ to obtain the generators $\dual^k$; set $\pred^{k+1} = (1 + h)\,\qq^k + \tstep\,\vel^k$; find the magnitude $s \ge 0$ such that $\pred^{k+1} - s\,\dual^k$ is realizable; set $\qq^{k+1} = \pred^{k+1} - s\,\dual^k$ and $\vel^{k+1} = (\qq^{k+1} - \qq^k)/\tstep$.

\paragraph{Projection (implicit).}
Maximize $\jointdual_\alpha$ with $\alpha = 1/h - 1$ and $\qq = \tcentroid = \tilde{\qq}^{k+1}$; set $\qq^{k+1}$ to the resulting centroids and $\vel^{k+1} = (\qq^{k+1} - \qq^k)/\tstep$.

\paragraph{Projection (explicit potential).}
Maximize $\jointdual_\alpha$ with $\qq = \qq^k$ and $\tcentroid = \pred^{k+1} - \qq^k/\alpha$, for any $\alpha > 0$ (\eg $\alpha = 1/h$, which gives $\tcentroid = \tilde{\qq}^{k+1}$); set $\qq^{k+1}$ to the resulting centroids and $\vel^{k+1} = (\qq^{k+1} - \qq^k)/\tstep$.

\paragraph{Penalty.}
Maximize $\jointdual_\alpha$ with $\qq = \tcentroid = \qq^k$ to obtain the centroids $\centroid^k$; set $\vel^{k+1} = \vel^k + (1 + \alpha)\,\tfrac{\tstep}{\stiff^2}(\centroid^k - \qq^k)$ and $\qq^{k+1} = \qq^k + \tstep\,\vel^{k+1}$.

\paragraph{\citet{Gallouet2018Lagrangian}.}
Penalty with $\alpha = 0$: maximize $\dualenergy$ with sites $\qq^k$ to obtain the centroids $\centroid^k$; set $\vel^{k+1} = \vel^k + \tfrac{\tstep}{\stiff^2}(\centroid^k - \qq^k)$ and $\qq^{k+1} = \qq^k + \tstep\,\vel^{k+1}$.

\paragraph{Power Particles \cite{deGoes2015PowerParticles}.}
Maximize $\dualenergy$ with sites $\qq^k$ to obtain the centroids $\centroid^k$; project the velocities, $\vel^{k+1} = \vel^k - \tstep\,\Mmat^{-1}\Gmat\,\press$ with $\tstep\,\Lmat\,\press = \Dmat\,\vel^k$; set $\qq^{k+1} = \centroid^k + \tstep\,\vel^{k+1}$.

\subsection*{Common skeleton of the update rules}

All schemes share the same skeleton: predict, correct along a normal of $\achievable$, and update the velocities.
They differ in where the normal is evaluated and how its magnitude is found, as listed in Table~\ref{tab:plan-skeleton}, with $\smash{\pred^{k+1} = (1+h)\,\qq^k + \tstep\,\vel^k}$ and $\smash{\tilde{\qq}^{k+1} = \qq^k + \tstep\,\vel^k}$.
\begin{table*}[t]
  \caption{[Plan] Common skeleton of the update rules.}
  \label{tab:plan-skeleton}
  \centering
  \footnotesize
  \setlength{\tabcolsep}{3pt}
  \begin{tabular}{@{}p{0.17\textwidth}p{0.19\textwidth}p{0.2\textwidth}p{0.19\textwidth}p{0.19\textwidth}@{}}
    \toprule
    scheme & prediction & correction direction & magnitude & solves per step \\
    \midrule
    SHAKE & $\pred^{k+1}$ & $\dual^k$, normal at $\qq^k$ & 1D search until realizable & one $\jointdual$, several $\jointdual_\alpha$ \\
    Projection (explicit) & $\pred^{k+1}$ & $\dual^{k+1}$, normal at $\qq^{k+1}$ & set by the projection & one $\jointdual_\alpha$ \\
    Projection (implicit) & $\tilde{\qq}^{k+1}/(1-h)$ & $\dual^{k+1}$, normal at $\qq^{k+1}$ & set by the projection & one $\jointdual_\alpha$ \\
    Power Particles & $\tilde{\qq}^{k+1}$, used as sites & $\tilde{\qq}^{k+1}$ (fixed) & $\to \infty$ (exposed point) & one $\dualenergy$, one pressure solve \\
    Penalty & none (explicit force) & $\qq^k - \centroid^k$ & stiffness $1 + \alpha$ & one $\jointdual_\alpha$ \\
    \citet{Gallouet2018Lagrangian} & none (explicit force) & $\qq^k - \centroid^k$ & stiffness $1$ & one $\dualenergy$ \\
    \bottomrule
  \end{tabular}
\end{table*}
SHAKE and the explicit projection differ in a single entry, the correction direction, which is the relaxation of \S\ref{sec:flow-projection}; the two projections differ only in the prediction, \ie in the treatment of the potential; Power Particles are the limit in which the prediction becomes the normal rather than the projected point; and the two penalty schemes have no correction step, since the constraint is replaced by a force.

\subsection*{SHAKE: computing the line search}

\paragraph{Why a single scalar.}
Within the affine subspace of \eqref{eq:partition-feasibility}, the realizable configurations form the relative boundary of $\achievable$ (minus its flat faces), a set of codimension one whose outward normal at a realizable $\qq$ is the generator vector $\dual$ \cite[Prop.~3.5]{Bourne2026ConvexOrder}.
The constraint thus behaves as SHAKE with a single bond: one scalar multiplier $s$, with the force along the normal $\dual^k$.
Unlike classical SHAKE, however, the constraint function is never written explicitly, and we can only test whether a point belongs to $\achievable$.

\paragraph{One time step.}
\begin{enumerate}
  \item \emph{Direction.} Maximize $\jointdual$ with sites and target centroids at $\qq^k$, \ie solve the inverse problem of \citet{Bourne2025Inverting}, to obtain the generators $\dual^k$ up to a homothety; center them as $\smash{\hat{\dual}^k = \dual^k - \tfrac{1}{\domainVol}\sum_i \pvol_i\,\dual^k_i}$, so that moving along $\hat{\dual}^k$ preserves \eqref{eq:partition-feasibility-centroid}.
  \item \emph{Prediction.} Set $\smash{\pred^{k+1} = (1 + h)\,\qq^k + \tstep\,\vel^k}$, shifted by the uniform translation that restores \eqref{eq:partition-feasibility-centroid}, which fixes the translation mode.
  \item \emph{Search.} Along the ray $\smash{\centroid(s) = \pred^{k+1} - s\,\hat{\dual}^k}$, $s \ge 0$, find the first $s$ at which $\centroid(s)$ enters $\achievable$:
    \begin{itemize}
      \item \emph{Membership oracle:} one maximization of $\jointdual_\alpha$, with sites $\qq^k$ and target centroids $\centroid(s) - \qq^k/\alpha$, returns the projection $\smash{P_\achievable(\centroid(s))}$ (\S\ref{sec:regularized-joint}), and thus the distance
      \begin{equation*}
        d(s) = \norm{\centroid(s) - P_\achievable(\centroid(s))}_{\pvol} .
      \end{equation*}
      \item \emph{Convexity:} $d$ is the distance to a convex set along a line, so it is convex in $s$ and vanishes exactly on an interval $[s_-, s_+]$; its slope is available from the same solve, $\smash{\tfrac{d}{ds}\tfrac12 d^2 = -\langle \centroid(s) - P_\achievable(\centroid(s)),\, \hat{\dual}^k \rangle_{\pvol}}$.
      \item \emph{Procedure:} first minimize $d$ over $s \ge 0$ by a safeguarded Newton or golden-section search; if $\min d > 0$, the ray misses $\achievable$ and the step has no solution, which is the failure mode of \S\ref{sec:flow-centroids}, now detected explicitly. Otherwise, bisect on $[0, s^\star]$, with any $s^\star$ such that $d(s^\star) = 0$, to locate the entry point $s_-$.
    \end{itemize}
  \item \emph{Update.} Set $\smash{\qq^{k+1} = \centroid(s_-)}$ and $\smash{\vel^{k+1} = (\qq^{k+1} - \qq^k)/\tstep}$.
\end{enumerate}

\paragraph{Cost and remarks.}
Each step costs one maximization of $\jointdual$ and about $\log(1/\text{tolerance})$ maximizations of $\jointdual_\alpha$, each warm-started from the previous trial.
\begin{itemize}
  \item \emph{Entry versus exit:} the ray crosses $\partial\achievable$ twice, at $s_-$ and $s_+$; SHAKE takes the root closest to the unconstrained motion, \ie $s_-$.
  \item \emph{Next direction:} at $\centroid(s_-) \in \partial\achievable$, the projection returns vanishing generators, so $\dual^{k+1}$ must come from a new maximization of $\jointdual$, which can be warm-started from the trials just outside, $s < s_-$.
  \item \emph{Flat faces:} if $\centroid(s_-)$ falls in the relative interior of a flat face of $\achievable$, it lies on the boundary but is not realizable.
  \item \emph{Relaxation:} the projection of \S\ref{sec:flow-projection} replaces the whole search by its single final projection, hence its single solve.
\end{itemize}

\paragraph{Explicit root-finding alternative.}
The constraint can also be written explicitly as $g(\qq) = 0$ with $\smash{g = \gamma_\achievable - 1}$, with the gauge function $\gamma_\achievable$ of $\achievable$ centered at $\smash{\domainCentroid\mathbf{1}}$, \ie the radial distance to its boundary.
Evaluating $\gamma_\achievable$ is exactly what the dropped ``$\jointdual$ on a hyperplane'' formulation computed (with $1 + r$ as the radial factor; see \texttt{gauge\_projection\_log.tex}), so classical SHAKE root-finding, \ie Newton on $g(\centroid(s)) = 0$, would apply directly.

\subsection*{An implicit symplectic integrator}

\paragraph{Why the projection schemes break symplecticity.}
Both projection schemes put the constraint normal at the \emph{new} configuration, $\smash{\dual^{k+1}}$, but in a single position update, so that the step is a non-invertible map: every prediction inside $\achievable$ is collapsed onto itself, and every prediction outside onto its boundary.
The discrete Euler--Lagrange equations of \S\ref{sec:flow-lagrangian}, in contrast, put the constraint normal at the \emph{middle} configuration $\qq^k$ of the two steps $(\qq^{k-1}, \qq^k)$ and $(\qq^k, \qq^{k+1})$, which is exactly SHAKE.
Any variational integrator obtained from a discrete Lagrangian on $\achievable$ inherits this property, so a symplectic scheme must keep the normal at $\qq^k$, \ie it must be a SHAKE-type scheme, possibly with a different quadrature of the potential.

\paragraph{A symplectic scheme with the normal at the new point.}
The way out is to split the step in two half-steps, as in RATTLE \cite{Andersen1983Rattle}, which is the symplectic integrator on the constraint manifold obtained from the symmetric discrete Lagrangian \cite[\S3.5.4]{marsden2001discrete}:
\begin{enumerate}
  \item \emph{Position step (SHAKE):}
  \begin{equation*}
    \vel^{k+1/2} = \vel^k + \tfrac{\tstep}{2\stiff^2}(\qq^k - \dual^k) , \qquad \qq^{k+1} = \qq^k + \tstep\,\vel^{k+1/2} ,
  \end{equation*}
  along the normal $\smash{\dual^k}$ at $\qq^k$, with its magnitude chosen so that $\qq^{k+1}$ is realizable, as in \S\ref{sec:flow-centroids};
  \item \emph{Velocity step:}
  \begin{equation*}
    \vel^{k+1} = \vel^{k+1/2} + \tfrac{\tstep}{2\stiff^2}(\qq^{k+1} - \dual^{k+1}) ,
  \end{equation*}
  along the normal $\smash{\dual^{k+1}}$ at $\qq^{k+1}$, with its magnitude chosen so that $\vel^{k+1}$ is tangent to the realizable configurations, \ie $\smash{\langle \vel^{k+1}, \dual^{k+1} \rangle_{\pvol} = 0}$ after removing the uniform translation.
\end{enumerate}
The second half-step is linear in the magnitude, and therefore solved in closed form; it puts a constraint normal at the new point, but only on the velocities, so the map remains invertible and symplectic.
Power Particles share with RATTLE the use of a separate velocity state, but not the velocity step itself: their pressure projection enforces $\Dmat\,\vel = 0$, \ie about $n$ linearized volume constraints at fixed weights, whereas the velocity step of RATTLE removes a single component along $\smash{\hat{\dual}^{k+1}}$, since the weights already enforce the volumes exactly. The pressure projection is instead the Schur-complement weight solve of the Hessian split below, \ie a first-order predictor for the exact weight solve that Power Particles run next.

\paragraph{Relation to the projection schemes.}
The projection schemes can be read as RATTLE with both half-steps merged into a single position projection: the position step uses the normal at the new point instead of the old one, and the velocity step is replaced by the difference quotient $\smash{(\qq^{k+1} - \qq^k)/\tstep}$, which does not remove the normal component of the velocity.
The latter is the source of the energy dissipation of the projection schemes (an inelastic impact), whereas RATTLE removes the normal component of the velocity explicitly and conserves a modified energy.
RATTLE thus requires the SHAKE position step, with its line search, but shares with the projection schemes the use of the normal at the new configuration.

\paragraph{An implicit symplectic variant.}
A fully implicit symplectic scheme is obtained with the implicit midpoint rule \cite[Example~2.3.2]{marsden2001discrete}, \ie by evaluating both the potential and the constraint at $\smash{\tfrac12(\qq^k + \qq^{k+1})}$.
Since this midpoint is generally not realizable, this variant requires the constraint at points of $\achievable$ that are not power diagrams, \eg through the gauge function $\gamma_\achievable$ of the dropped site-free formulation, and it is substantially more costly than RATTLE.
We therefore propose RATTLE as the symplectic counterpart of the projection schemes, and keep the projections as cheap, dissipative alternatives.

\paragraph{Consequences for Section~\ref{sec:flow-lagrangian}.}
Since the motivation of \S\ref{sec:flow} is variational, a possible reorganization is to present SHAKE and RATTLE as the two variational integrators of the constrained Lagrangian, the penalty scheme as the variational integrator of the softened Lagrangian, and the projection schemes as cheaper but non-variational relaxations; Power Particles then appear as a relaxation of the projection family, whose pressure projection corresponds to the weight--shift split of the Newton step on $\jointdual$.

\paragraph{To check.}
\begin{itemize}
  \item Whether the realizable configurations form a smooth manifold near the current configuration, which RATTLE requires; Bourne et al.\ show that they are open in the relative boundary of $\achievable$ \cite[Thm.~3.10]{Bourne2026ConvexOrder}, but smoothness is not established.
  \item Whether the velocity step of RATTLE, which removes one normal component, dissipates less than the pressure projection of Power Particles, which removes about $n$ components.
\end{itemize}

\subsection*{Line-constrained $\jointdual$: success and failure}

\paragraph{Problem.}
The position step of SHAKE and RATTLE can be computed in a single solve by constraining the shifts to a hyperplane, as in the dropped site-free formulation (\texttt{gauge\_projection\_log.tex}):
\begin{equation*}
  \max_{\shift, \weight}\ \jointdual(\shift, \weight)
  \quad \text{s.t.} \quad
  \langle \shift, \mathbf{d} \rangle_{\pvol} = 0 ,
\end{equation*}
with sites $\qq^k$, target centroids at the prediction $\pred^{k+1}$, and the centered normal $\smash{\mathbf{d} = \hat{\dual}^k}$ at $\qq^k$.

\paragraph{Reduction to one dimension.}
The multiplier $r$ of the constraint moves the targets along the line $\smash{L = \{\pred^{k+1} + r\,\mathbf{d}\}}$.
For fixed $r$, maximizing over the shifts and weights returns the transport cost $\smash{\int_\domain\norm{\xx}^2 d\xx + \norm{\qq^k}^2_{\pvol} - 2\langle \qq^k, \pred^{k+1} + r\,\mathbf{d} \rangle_{\pvol}}$ if $\smash{\pred^{k+1} + r\,\mathbf{d} \in \achievable}$, and $+\infty$ otherwise (Appendix~\ref{sec:appendix-existence}).
The problem thus reduces to minimizing $\smash{-2r\,\langle \qq^k, \mathbf{d} \rangle_{\pvol}}$ over the interval $\smash{I = \{ r : \pred^{k+1} + r\,\mathbf{d} \in \achievable \}}$.

\paragraph{Endpoint selection.}
Since $\hat{\dual}^k$ is an outward normal at $\qq^k$, and the configuration $\smash{\domainCentroid\mathbf{1}}$ lies in the relative interior of $\achievable$, we have $\smash{\langle \qq^k, \hat{\dual}^k \rangle_{\pvol} > 0}$, so the solution is $r^\star = \max I$, \ie the outermost point of $\achievable$ on $L$.
For a prediction outside $\achievable$, this is the first entry point coming from the prediction, \ie the near root of SHAKE; for a prediction inside $\achievable$, $r^\star > 0$ pushes it outward to the boundary.
The constraint is thus bilateral, and the multiplier may take either sign, contrary to the search over $s \ge 0$ described above.

\paragraph{Success and failure.}
The solve succeeds whenever
\begin{itemize}
  \item[(S1)] the line $L$ meets $\achievable$, \ie $I \neq \emptyset$;
  \item[(S2)] $\smash{\langle \qq^k, \mathbf{d} \rangle_{\pvol} \neq 0}$, which always holds, as shown above;
  \item[(S3)] the endpoint $r^\star$ is an exposed point of $\achievable$, rather than a point in the relative interior of a flat face;
  \item[(S4)] the Newton iterates keep the generators distinct.
\end{itemize}
Condition (S1) fails only when the tangential motion carries the prediction beyond the silhouette of $\achievable$ seen along $\hat{\dual}^k$; since the projection of $\qq^k$ along $\hat{\dual}^k$ lies generically inside this silhouette, (S1) holds for small $\tstep$, and a tangent velocity, as enforced by RATTLE, widens the admissible range; when (S1) fails, $\jointdual$ is unbounded, which is detectable.
Condition (S3) holds for small steps from a realizable $\qq^k$, since the exposed points form an open subset of the relative boundary \cite[Thm.~3.10]{Bourne2026ConvexOrder}, but flat faces cover a positive-measure part of the boundary and may be reached by large steps.
The hard-constrained solve is thus reliable for small steps, and fails detectably for large ones.

\paragraph{Softened fallback.}
The softened dual $\jointdual_\alpha$ alone cannot replace this solve, since it places the normal at the new point and thus changes the integrator into the projection scheme.
Keeping the linear constraint, however, the problem
\begin{equation*}
  \max_{\shift, \weight}\ \jointdual_\alpha(\shift, \weight)
  \quad \text{s.t.} \quad
  \langle \shift, \hat{\dual}^k \rangle_{\pvol} = 0
\end{equation*}
is, by the same duality, equivalent to
\begin{equation*}
  \min_{\centroid \in \achievable}\ \alpha\,\mathrm{dist}^2_{\pvol}(\centroid, L) - 2\,\langle \qq^k, \centroid \rangle_{\pvol} ,
\end{equation*}
which always has a minimizer, since $\achievable$ is compact and the objective is continuous.
If $L$ meets $\achievable$, the minimizer tends to the hard-constrained solution as $\alpha \to \infty$, with an overshoot of order $1/\alpha$ beyond the endpoint; otherwise, it returns the point of $\achievable$ closest to $L$, tilted toward $\qq^k$, so the scheme degrades gracefully instead of failing.
The remaining caveat is the distinctness of the generators: the bound \eqref{eq:regularized-distinct} requires small $\alpha$, whereas closeness to SHAKE requires large $\alpha$, so distinctness is only guaranteed in the hard-constrained limit, under (S3).

\paragraph{To check.}
\begin{itemize}
  \item The silhouette argument for (S1), which assumes $\achievable$ to be smooth and strictly convex near $\qq^k$.
  \item The rate of the $O(1/\alpha)$ overshoot of the softened fallback.
  \item The uniqueness of the softened minimizer, whose objective is convex but not strictly convex in $\centroid$.
\end{itemize}

\subsection*{Implicit treatment of the penalty}

\paragraph{Classification.}
The only potential of our formulations is the stiff penalty $\smash{\tfrac{1}{2\stiff^2}\cpot_\alpha}$, whose stiffness grows as $\smash{(1+\alpha)/\stiff^2}$, both in the limit $\stiff \to 0$ of \citet{Gallouet2018Lagrangian} and in the constrained limit $\alpha \to \infty$, where it becomes the realizability constraint.
Our schemes then differ in how they treat this potential and this constraint, as listed in Table~\ref{tab:plan-imex}, which reveals one missing combination: a symplectic, implicit treatment of the potential.

\begin{table*}[t]
  \caption{[Plan] Treatment of the potential and of the constraint in each scheme.}
  \label{tab:plan-imex}
  \centering
  \small
  \begin{tabular}{@{}llll@{}}
    \toprule
    scheme & potential & constraint & symplectic \\
    \midrule
    \citet{Gallouet2018Lagrangian}, penalty & explicit & --- & yes (step limit $\tstep \lesssim \stiff$) \\
    SHAKE / RATTLE & explicit & implicit, normal at $\qq^k$ & yes \\
    Projection (explicit) & explicit & implicit, normal at $\qq^{k+1}$ & no \\
    Projection (implicit) & implicit Euler & implicit, normal at $\qq^{k+1}$ & no \\
    Implicit midpoint (proposed) & implicit midpoint & --- (penalty) & yes \\
    \bottomrule
  \end{tabular}
\end{table*}

\paragraph{Implicit midpoint.}
Evaluating the penalty at the midpoint of the time step, as in the implicit midpoint rule of \S\ref{sec:flow-lagrangian} \cite[Example~2.3.2]{marsden2001discrete}, gives the discrete Lagrangian
\begin{equation*}
  \lag_d(\qq^k, \qq^{k+1}) = \frac{1}{2\tstep} \norm{\qq^{k+1} - \qq^k}^2_{\pvol} - \frac{\tstep}{2\stiff^2}\, \cpot_\alpha\Bigl( \frac{\qq^k + \qq^{k+1}}{2} \Bigr) ,
\end{equation*}
\ie the penalty scheme of \S\ref{sec:flow-softened} with the midpoint quadrature mentioned in \S\ref{sec:flow-lagrangian} instead of the evaluation at $\qq^k$.
Since it derives from a discrete Lagrangian, this integrator is variational, hence symplectic, and it keeps the motivation of \S\ref{sec:flow-lagrangian}.
It is stable for any $\alpha$ and $\stiff$, so $\alpha$ can be taken large, which keeps the particles close to the centroids, and $\stiff$ small, without the step limit of the explicit penalty.
Each step requires a nonlinear solve for $\qq^{k+1}$, in which the gradient of $\cpot_\alpha$ at the midpoint is $\smash{2\,\cellvol\,(\dual - \centroid)}$ by \eqref{eq:softened-site-gradient}, \ie one maximization of $\jointdual_\alpha$ with sites and target centroids at the midpoint; a fixed-point or Newton iteration, warm-started across iterations, thus costs a few solves per step, which is comparable to SHAKE.
As $\alpha \to \infty$, however, the implicit midpoint rule remains stable but does not reproduce the constrained dynamics exactly, since it aliases rather than removes the fast normal oscillations, a known effect for highly oscillatory problems; it therefore tracks the centroid dynamics only approximately.

\paragraph{Implicit Euler on the penalty of \citet{Gallouet2018Lagrangian}.}
For $\alpha = 0$ and an implicit Euler treatment of the potential, with $h = \tstep^2/\stiff^2$ and the drift $\smash{\tilde{\qq}^{k+1} = \qq^k + \tstep\,\vel^k}$, the update is the fixed point
\begin{equation*}
  (1 + h)\,\qq^{k+1} = \tilde{\qq}^{k+1} + h\,\centroid(\qq^{k+1}) ,
\end{equation*}
where $\centroid(\cdot)$ denotes the centroids of the volume-constrained power diagram with the given sites, \ie one maximization of $\dualenergy$.
On partitions with volumes $\pvol$, the potential reads
\begin{equation*}
  \cpot_0(\qq) = \text{const} + \norm{\qq}^2_{\pvol} - 2\,\sigma_\achievable(\qq) ,
\end{equation*}
with the support function $\sigma_\achievable$ of $\achievable$, so the implicit step minimizes a difference of convex functions.
Eliminating $\qq$, it is equivalent to
\begin{equation*}
  \max_{\centroid \in \achievable}\ \norm{\tilde{\qq}^{k+1} + h\,\centroid}^2_{\pvol} ,
\end{equation*}
whose local maxima, as maxima of a convex function over a convex set, are extreme points of $\achievable$, \ie realizable centroids.
The natural solver is the DC algorithm, whose monotone iteration
\begin{equation*}
  \qq \leftarrow \frac{\tilde{\qq}^{k+1} + h\,\centroid(\qq)}{1 + h}
\end{equation*}
costs one maximization of $\dualenergy$ per iteration, \ie a damped Lloyd step.

\paragraph{Power Particles as one DC iteration.}
Starting the DC algorithm at $\smash{\qq = \tilde{\qq}^{k+1}}$ and taking the stiff limit $h \to \infty$, \ie $\stiff \to 0$, a single iteration gives $\smash{\qq^{k+1} = \centroid(\tilde{\qq}^{k+1})}$, the centroids of the power diagram generated by the drifted points, which is exactly the position step of Power Particles, up to the cyclic rotation of \S\ref{sec:flow-pp}.
This agrees with the identification of Power Particles with the projection for $\alpha \to 0$ (Table~\ref{tab:plan-family}), but now with a dynamical meaning: the position step of Power Particles is one Lloyd-type iteration of implicit Euler on the penalty of \citet{Gallouet2018Lagrangian}, in the stiff limit, and their pressure projection is an additional velocity correction, related to the weight--shift split of the Newton step on $\jointdual$.

\paragraph{Three levels of the same idea.}
\begin{enumerate}
  \item \emph{Power Particles:} one DC iteration in the stiff limit, followed by a pressure projection; cheap and dissipative.
  \item \emph{Implicit Euler on the penalty, iterated to convergence:} a few $\dualenergy$ solves per step, with fixed points closer to realizable configurations as $h$ grows; still dissipative and not symplectic.
  \item \emph{Implicit midpoint on $\cpot_\alpha$:} symplectic and unconditionally stable in the stiff part, with a few $\jointdual_\alpha$ solves per step, and centroids tracked up to $O(1/\alpha)$.
\end{enumerate}
The last level matches best the goal of a variational, implicit integrator for the centroids of a moving power diagram with prescribed volumes, while the first two explain Power Particles as a cheap truncation of it.

\paragraph{To check.}
\begin{itemize}
  \item The identity $\smash{\cpot_0 = \text{const} + \norm{\qq}^2_{\pvol} - 2\,\sigma_\achievable(\qq)}$, which follows from the expansion of the transport cost used for \eqref{eq:regularized-projection}.
  \item Whether the one-step identification with Power Particles survives the cyclic rotation of their algorithm.
  \item The convergence and cost of the fixed-point solver of the midpoint scheme.
  \item The aliasing of the normal oscillations for large $\alpha$, on a 2D example.
\end{itemize}

\subsection*{Comparison (draft of the discussion table)}
Table~\ref{tab:plan-comparison} compares the properties of all schemes.

\begin{table*}[t]
  \caption{[Plan] Comparison of the integrators (draft of the discussion table).}
  \label{tab:plan-comparison}
  \centering
  \footnotesize
  \setlength{\tabcolsep}{3pt}
  \begin{tabular}{@{}p{0.13\textwidth}p{0.105\textwidth}p{0.105\textwidth}p{0.105\textwidth}p{0.105\textwidth}p{0.105\textwidth}p{0.105\textwidth}p{0.105\textwidth}@{}}
    \toprule
    & SHAKE & Softened SHAKE & Proj.\ (expl.) & Proj.\ (impl.) & Penalty & \citet{Gallouet2018Lagrangian} & Power Particles \\
    \midrule
    particles are & centroids & centroids & centroids & centroids & sites & sites & sites \\
    dual & $\jointdual$ on a line & $\jointdual_\alpha$ on a line & $\jointdual_\alpha$ & $\jointdual_\alpha$ & $\jointdual_\alpha$ & $\dualenergy$ & $\dualenergy$ \\
    solves per step & one (bordered) & one (bordered) & one & one & one & one & one + linear solve \\
    exact centroids & yes & $O(1/\alpha)$ & yes$^*$ & yes$^*$ & $O(1/\alpha)$ & no & no \\
    exact volumes & yes & yes & yes & yes & yes & yes & yes \\
    symplectic & yes & as $\alpha \to \infty$ & no & no & yes & yes & no \\
    potential & explicit & explicit & explicit & implicit & explicit & explicit & none \\
    time-step limit & small $\tstep$ & none & none & $\tstep < \stiff$ & ? & $\tstep \le \stiff$ & CFL \\
    failure modes & $L \cap \achievable = \emptyset$, flat faces & merging generators & inward drift, flat faces & inward drift, flat faces & merging generators & none & none \\
    energy & preserved (modified) & approx.\ preserved & non-increasing$^\dagger$ & non-increasing$^\dagger$ & preserved (modified) & preserved (modified) & ? \\
    \bottomrule
  \end{tabular}
  \\[1mm]
  {\footnotesize $^*$achievable, and realizable unless the projection falls inside $\achievable$ or in a flat face. $^\dagger$without potential.}
\end{table*}

\subsection*{Power Particles as an inexact Newton step on $\jointdual$}

\paragraph{Newton system.}
We denote the negated Hessian blocks of $\jointdual$ by $\smash{\Kmat = -\partial^2_{\shift\shift}\jointdual}$ (shift--shift), $\smash{\wlap = -\partial^2_{\weight\weight}\jointdual}$ (the volume Laplacian of the shifted sites), and $\smash{\mathrm{B} = \partial^2_{\shift\weight}\jointdual = -\partial\cellvol/\partial\shift}$ (mixed, \eqref{eq:joint-hessian}).
Since $\jointdual$ is concave, one Newton step solves
\begin{equation*}
  \begin{bmatrix} \Kmat & -\mathrm{B} \\ -\mathrm{B}\tr & \wlap \end{bmatrix}
  \begin{bmatrix} \delta\shift \\ \delta\weight \end{bmatrix}
  =
  \begin{bmatrix} \partial_{\shift}\jointdual \\ \partial_{\weight}\jointdual \end{bmatrix} .
\end{equation*}
Eliminating the shifts splits this step into two stages:
\begin{align*}
  \bigl( \wlap - \mathrm{B}\tr\Kmat^{-1}\mathrm{B} \bigr)\, \delta\weight &= \partial_{\weight}\jointdual + \mathrm{B}\tr\Kmat^{-1}\partial_{\shift}\jointdual , \\
  \delta\shift &= \Kmat^{-1}\bigl( \partial_{\shift}\jointdual + \mathrm{B}\,\delta\weight \bigr) ,
\end{align*}
\ie a weight (pressure) solve followed by a shift (velocity) correction.
The reduced operator $\smash{\wlap - \mathrm{B}\tr\Kmat^{-1}\mathrm{B}}$ is the Schur complement of the shift block, and the back-substitution corrects the shifts along $\smash{\Kmat^{-1}\mathrm{B}\,\delta\weight}$, \ie along a discrete gradient of the weight update.

\paragraph{Pressure projection of Power Particles.}
Power Particles project the velocities by minimizing $\smash{\tfrac12\norm{\vel^+ - \vel}^2_{\Mmat}}$ subject to $\smash{\Dmat\,\vel^+ = 0}$, which is solved by $\smash{\tstep\,\Lmat\,\press = \Dmat\,\vel}$ and $\smash{\vel^+ = \vel - \tstep\,\Mmat^{-1}\Gmat\,\press}$, with $\smash{\Gmat = -\Dmat\tr}$ and $\smash{\Lmat = \Dmat\,\Mmat^{-1}\Gmat}$ \cite{deGoes2015PowerParticles}.
The two stages above match this projection term by term, as listed in Table~\ref{tab:plan-pp-newton}.
\begin{table*}[t]
  \caption{[Plan] Newton on $\jointdual$ versus the pressure projection of Power Particles.}
  \label{tab:plan-pp-newton}
  \centering
  \small
  \begin{tabular}{@{}ll@{}}
    \toprule
    Newton on $\jointdual$ & Power Particles \\
    \midrule
    $\mathrm{B} = -\partial\cellvol/\partial\shift$ & $-\Dmat\tr$ (divergence with respect to the sites, at fixed weights) \\
    $\Kmat$ & $2\,\Mmat$ (lumped mass) \\
    $\mathrm{B}\tr\Kmat^{-1}\mathrm{B}$ & $\tfrac12\,\Dmat\,\Mmat^{-1}\Dmat\tr = -\tfrac12\,\Lmat$ (pressure Laplacian) \\
    weight update $\delta\weight$ & pressure $\press$ \\
    shift correction $\Kmat^{-1}\mathrm{B}\,\delta\weight$ & velocity correction $\Mmat^{-1}\Gmat\,\press$ \\
    \bottomrule
  \end{tabular}
\end{table*}
Two observations support this dictionary.
First, by Appendix~\ref{sec:appendix-centroid}, each diagonal block of $\Kmat$ is $\smash{2\,\cellvol_i\,\Id}$ plus facet integrals, since the shifted weight $\combweight_i$ of \eqref{eq:shifted-site} has second derivative $\smash{\partial^2_{\shift_i}\combweight_i = 2\,\Id}$, so the lumped mass matrix of Power Particles is $\Kmat/2$ without the facet terms; with the softening of \S\ref{sec:regularized-joint}, $\jointdual_\alpha$ adds $\smash{\tfrac{2}{\alpha}\,\Mmat}$, which dominates as $\alpha \to 0$, \ie at the end of the projection family where Power Particles sit.
Second, by \eqref{eq:appendix-jac}, the shifted weight varies as $\smash{\partial_{\shift_i}\combweight_i = 2\,(\dual_i - \tcentroid_i)}$, which vanishes at zero shifts when the targets coincide with the sites, as in the projection with sites and targets at the drift; there, $\smash{\partial\cellvol/\partial\shift = \partial\cellvol/\partial\dual}$ is exactly the divergence $\Dmat$ of Power Particles.

\paragraph{Three simplifications.}
Starting from this dictionary, Power Particles depart from one exact Newton step in three ways:
\begin{enumerate}
  \item \emph{Weights first, to convergence.} They solve the weight block alone (Newton on $\dualenergy$) before the pressure projection, so that $\smash{\partial_{\weight}\jointdual = 0}$ when the projection starts, and build $\Dmat$ and $\Lmat$ on the power diagram of the converged weights; this is a block Gauss--Seidel ordering, weights then shifts.
  \item \emph{Lumped mass.} They replace $\Kmat/2$ by the diagonal $\Mmat$, dropping the facet terms.
  \item \emph{Single step at the velocity level.} They apply one Schur-complement step to the velocities rather than iterating the coupled system on the positions, and discard the remaining source $\smash{\partial_{\shift}\jointdual}$ together with the $\smash{O(\tstep^2)}$ position correction.
\end{enumerate}
In this reading, there are three related variants of one coupled step: the exact Newton step (Schur complement, weights coupled to shifts), the shifts-then-weights block Gauss--Seidel ordering, and the weights-then-shifts ordering with a lumped shift block, which is Power Particles.

\paragraph{To check.}
\begin{itemize}
  \item The constants and signs between $\delta\weight$ and $\tstep\,\press$, against \cite[Eqs.~(4) and~(7)]{deGoes2015PowerParticles}.
  \item The facet-term form of $\Kmat$ from \eqref{eq:centroid-hessian} and Appendix~\ref{sec:appendix-centroid}.
  \item The first-order equivalence between the position correction of our step and the velocity projection of Power Particles, through $\smash{\vel^{k+1} = (\qq^{k+1} - \qq^k)/\tstep}$.
\end{itemize}

\subsection*{Open decisions}

\begin{enumerate}
  \item \textbf{Presentation of the update rules:} braced equations as in \eqref{eq:flow-shake} for each scheme (cheapest, consistent with the rest of \S\ref{sec:flow}), one \texttt{algorithm} block per scheme (clearer for implementers, but longer), or a single figure with side-by-side pseudocode (visual comparison, but tight in two columns).
  \item \textbf{Explicit-potential projection:} keep it as a remark in \S\ref{sec:flow-projection}, or drop it to present a single projection scheme.
  \item \textbf{Table entries marked ``?''}: the time-step restriction of the penalty scheme, and the energy behavior of Power Particles, need to be checked.
\end{enumerate}

\end{document}
